\section{Phase 2 --- Chọn cách hỏi LLM}

Phase 2 nhận \NumPoneRerankTopN{} đoạn văn Phase 1 đưa sang và trả lời một câu
hỏi duy nhất: ra lệnh cho Gemini thế nào để nó trả lời đúng mà không bịa. Cần
nói rõ phạm vi --- \textbf{không có huấn luyện lại hay fine-tune nào cả}. Chỉ
hai thứ được thay đổi: nội dung system prompt, và số đoạn context đưa vào.

\subsection{Baseline, và chẩn đoán từ ba con số}

Chạy prompt gốc P0 với \NumWinnerDepth{} đoạn context trên \NumNDev{} câu
development:

\begin{center}\small
\begin{tabular}{@{}c c c c@{}}
\toprule
\textbf{Answer Correctness} & \textbf{Exact Match} & \textbf{Token F1} & \textbf{Faithfulness} \\
\midrule
\NumPzeroAC & \NumPzeroEM & \NumPzeroFone & \NumPzeroFaith \\
\bottomrule
\end{tabular}
\end{center}

Ba con số này đọc cùng nhau cho một chẩn đoán duy nhất. Faithfulness
\NumPzeroFaith{} nghĩa là gần như mọi khẳng định đều bám vào bằng chứng đã lấy
được --- hệ thống \emph{không} bịa. Nhưng Exact Match \NumPzeroEM{} nghĩa là
không một câu trả lời nào trùng khớp đáp án chuẩn.

\keypoint{\textbf{Hệ thống không bịa. Nó chỉ nói dài.} Model biết đáp án nhưng
gói nó trong một đoạn văn, trong khi đáp án chuẩn của NewsQA là một span rất
ngắn: một cái tên, một mốc thời gian, một con số. Vấn đề nằm ở \emph{hình dạng}
câu trả lời, không phải ở độ trung thực của nó.}

Chẩn đoán này đến từ một audit tay \textbf{30 câu} có Answer Correctness thấp
nhất, phân loại theo nguyên nhân. Đáng chú ý: người duyệt thấy \textbf{23/30}
câu điểm thấp thực ra \emph{đúng về mặt ngữ nghĩa} --- điểm thấp vì cách diễn
đạt, hoặc vì nhãn có vấn đề.

\subsection{Bốn prompt, mỗi cái nhắm một nhóm lỗi đã đếm được}

Prompt không được viết theo trực giác mà theo bảng phân loại lỗi từ audit. Nội
dung nguyên văn cả bốn prompt nằm ở \texttt{docs/prompts.md}.

\begin{center}\small
\begin{tabularx}{0.94\textwidth}{@{}p{0.9cm} p{3.4cm} L p{1.2cm}@{}}
\toprule
\textbf{ID} & \textbf{Ý tưởng} & \textbf{Nhắm nhóm lỗi} & \textbf{Cỡ nhóm} \\
\midrule
P0 & Prompt RAG có citation & \emph{control} & --- \\
P1 & Siết grounding tường minh & Hallucination & $\sim$0 \\
P2 & Trả lời ngắn, đúng answer type & Đúng nhưng quá dài & 11/30 \\
P3 & Chọn đúng đoạn trước khi trả lời & Trộn bài, trộn mốc thời gian & 3/30 \\
\bottomrule
\end{tabularx}
\end{center}

\paragraph{P1 thất bại, và đó là bài học đắt nhất.}
P1 siết ràng buộc bằng chứng để tăng faithfulness. Nhưng baseline đã đạt
\NumPzeroFaith{} --- không còn dư địa. P1 nhắm vào một vấn đề không tồn tại và
\emph{mất} correctness khi làm vậy. Trực giác ``siết chặt hơn thì tốt hơn'' là
sai khi chưa đo xem nhóm lỗi ấy có bao nhiêu ca.

\paragraph{P2 thắng vì nó khớp với hình dạng của nhãn.}
P2 yêu cầu một câu trả lời ngắn chứa đúng loại thông tin được hỏi. Đây cũng là
lý do Answer Relevancy \emph{giảm} --- metric ấy thưởng câu trả lời đầy đủ ngữ
cảnh, còn dữ liệu này thưởng câu trả lời cụt. Sự giảm đó là hệ quả có chủ đích,
không phải suy thoái.

\subsection{\texttt{context\_depth} --- thứ được tinh chỉnh song song}

\texttt{context\_depth} là số đoạn được dán vào khối context, lấy từ đầu danh
sách đã xếp hạng. Trong mã nguồn nó đúng là một phép cắt danh sách:

\begin{center}
\texttt{generation\_chunks = ranked\_chunks[:context\_depth]}
\end{center}

Không truy xuất lại, không xếp hạng lại --- nên chênh lệch giữa các mức depth
quy hết được cho tầng sinh. Model nhận khối context đánh số \texttt{[1] [2]
[3]\ldots}, và chính các chỉ số này là số nó dùng khi trích dẫn.

\begin{center}\small
\begin{tabularx}{0.94\textwidth}{@{}p{1.4cm} p{2.6cm} L@{}}
\toprule
\textbf{depth} & \textbf{Model thấy} & \textbf{Đánh đổi} \\
\midrule
1 & \texttt{[1]} & Ít token nhất --- nhưng \NumLostAtDepthOne{}/\NumNDev{} câu mất sạch bằng chứng \\
3 & \texttt{[1] [2] [3]} & Cân bằng --- nhưng \NumLostAtDepthThree{}/\NumNDev{} câu mất sạch bằng chứng \\
5 & \texttt{[1]}\ldots\texttt{[5]} & Không mất bằng chứng --- tốn token nhất \\
\bottomrule
\end{tabularx}
\end{center}

Cột bên phải không đến từ ước lượng. Nó đọc thẳng ra từ đường cong Hit@k mà
Phase~1 đã đo: Hit@1 \NumPoneHitOne{}, Hit@3 \NumPoneHitThree{}, Hit@5
\NumPoneHitFive{}. Đây là chỗ Phase~1 định giá sẵn một quyết định của Phase~2
mà không tốn thêm một lệnh gọi API nào.

\subsection{Quyết định winner: guardrail chạy trước, điểm số chạy sau}

Bốn điều kiện được khóa \emph{trước khi nhìn thấy bất kỳ kết quả nào}:
Faithfulness không tụt quá $0{,}02$; Citation F1 không tụt quá $0{,}01$;
Citation Validity không tụt quá $0{,}01$; coverage $\geq 95\%$. Chỉ những cấu
hình qua cả bốn mới được xét, và trong số đó mới chọn Answer Correctness cao
nhất.

\begin{center}\small
\begin{tabularx}{0.94\textwidth}{@{}L p{2cm} p{2.4cm} p{2.4cm}@{}}
\toprule
 & \textbf{P0} & \textbf{P2-depth3} & \textbf{P2-depth5} \\
\midrule
Answer Correctness & \NumPzeroAC & \emph{\NumPtwodThreeAC} & \NumPtwodFiveAC \\
Exact Match & \NumPzeroEM & \emph{\NumPtwodThreeDeltaEM} & \NumPtwodFiveDeltaEM \\
\midrule
Faithfulness $\Delta$ & --- & \NumPtwodThreeGuardFaith{} (\fail{FAIL}) & \NumPtwodFiveGuardFaith{} (\pass{PASS}) \\
Citation F1 $\Delta$ & --- & \NumPtwodThreeGuardCitFone{} (\pass{PASS}) & \NumPtwodFiveGuardCitFone{} (\pass{PASS}) \\
Citation Validity $\Delta$ & --- & \NumPtwodThreeGuardCitVal{} (\fail{FAIL}) & \NumPtwodFiveGuardCitVal{} (\pass{PASS}) \\
\midrule
\textbf{Phán quyết} & --- & \fail{LOẠI} & \win{CHỌN} \\
\bottomrule
\end{tabularx}
\end{center}

P2-depth3 đạt Answer Correctness \emph{cao hơn} (\NumPtwodThreeAC{} so với
\NumPtwodFiveAC{}) nhưng trượt 2 trong 4 guardrail. Kiểm định theo cặp xác nhận
cả hai mức tụt đều là thật, không phải nhiễu đo: Faithfulness
\NumPtwodThreeDeltaFaith{} với CI95 \NumPtwodThreeDeltaFaithCI{} --- không chứa 0.

\keypoint{Nới ngưỡng ở đúng thời điểm này --- sau khi đã thấy rằng cấu hình bị
loại có điểm cao hơn --- chính là hành vi mà việc đăng ký trước sinh ra để
ngăn. Bộ guardrail đã làm đúng việc của nó.}

\subsection{Winner: \NumWinner{}}

\begin{center}\small
\begin{tabularx}{0.94\textwidth}{@{}L p{2.2cm} p{3.4cm}@{}}
\toprule
So với baseline P0 (\NumNDev{} câu development) & \textbf{$\Delta$} & \textbf{CI95} \\
\midrule
Answer Correctness & \NumPtwodFiveDeltaAC & \NumPtwodFiveDeltaACCI \\
Exact Match & \NumPtwodFiveDeltaEM & \NumPtwodFiveDeltaEMCI \\
Token F1 & \NumPtwodFiveDeltaFone & \NumPtwodFiveDeltaFoneCI \\
Faithfulness & \NumPtwodFiveDeltaFaith & \NumPtwodFiveDeltaFaithCI \\
Citation F1 & \NumPtwodFiveDeltaCitFone & \NumPtwodFiveDeltaCitFoneCI \\
\bottomrule
\end{tabularx}
\end{center}

Correctness tăng \NumPtwodFiveDeltaAC{} với khoảng tin cậy không chứa 0.
Faithfulness \NumPtwodFiveDeltaFaith{} với CI95 \NumPtwodFiveDeltaFaithCI{}
--- \emph{không tách được khỏi 0}, và đây là kết quả tốt: cải thiện correctness
mà không đánh đổi grounding. Cái giá phải trả là chi phí input token cao hơn
depth~3 khoảng 60\%.

\subsection{Hai lát cắt kiểm tra}

\paragraph{Phân tầng theo truy xuất.}
Truy xuất đưa được bằng chứng ra trước mặt model ở \NumNHit{}/\NumNDev{} câu;
\NumNMiss{} câu còn lại thì không.

\begin{center}\small
\begin{tabularx}{0.94\textwidth}{@{}L p{1.2cm} p{1.8cm} p{2cm} p{2.4cm}@{}}
\toprule
Answer Correctness & \textbf{n} & \textbf{P0} & \textbf{P2-depth3} & \textbf{P2-depth5} \\
\midrule
\texttt{gold\_in\_top5} & \NumNHit & \NumPzeroHitAC & \NumPtwodThreeHitAC & \NumPtwodFiveHitAC \\
\texttt{gold\_not\_in\_top5} & \NumNMiss & \NumPzeroMissAC & \NumPtwodThreeMissAC & \NumPtwodFiveMissAC \\
\bottomrule
\end{tabularx}
\end{center}

Trên \NumNMiss{} câu không có gold trong context, đổi prompt chỉ nhích Answer
Correctness từ \NumPzeroMissAC{} lên \NumPtwodFiveMissAC{} --- và phần lớn mức
đó là điểm ngữ nghĩa cho một câu từ chối đúng cách, không phải câu trả lời
đúng. Trên \NumNHit{} câu có gold, cùng thay đổi prompt ấy tăng hơn $0{,}10$.

\paragraph{Macro theo bài báo.}
Trung bình theo câu để một bài đóng góp nhiều câu lấn át bài đóng góp ít câu.
Đọc theo macro (trung bình theo bài trước, rồi trung bình theo \NumNArticles{}
bài): P0 \NumPzeroACMacro{}, P2-depth3 \NumPtwodThreeACMacro{}, P2-depth5
\NumPtwodFiveACMacro{}. Cả ba mức chênh so với micro đều dưới $0{,}004$ ---
không bài báo nào chi phối kết quả, và winner không đổi dù đọc theo cách nào.
